\Qhead{2.4}{Волны в периодических структурах. Связанные волны}{Волны в периодических структурах. Механические цепочки, акустические и оптические фононы. Полосы пропускания и непрозрачности. Электрические цепочки, сплошная среда со слабыми периодическими неоднородностями. Связанные волны.}
\Subfile{2.4.1}{sections/02_waves/q04/s01}
\Subfile{2.4.2}{sections/02_waves/q04/s02}

\recap{Периодичность заменяет непрерывную симметрию дискретной. Дисперсионные ветви и зоны определяются ячейкой; брэгговская связь встречных волн открывает запрещённую полосу.}
\board{Одноатомная дисперсия; акустическая и оптическая ветви; условие Брэгга и $q^2=\delta^2-\kappa^2$.}
\oral{Почему в запрещённой полосе непоглощающей решётки энергия не обязана превращаться в тепло?}
\bait{Можно оставить континуальный предел: он правильно описывает звук при $ka\ll1$, но скрывает границу зоны и брэгговскую щель.}
\sources{Белокопытов~\cite{belokopytov1}, разделы о цепочках и слабонеоднородных системах; Цысарь~\cite{tsysar1}, лекция~10. Дисперсии проверены подстановкой в дискретные уравнения.}
